\subsection{Concentration of exact maximizers for varying densities}
In the proof of \cref{thm:constant-mass-entire-determinant-profile}, each current $T_r$ attains the sharp bound \eqref{eq:constant-mass-determinant-dual} for its own function $r$. Such exact maximizers, taken for varying $r$, need not have uniformly integrable critical powers $k_T^{n/(n-1)}$, and their weak limits can have zero determinant measure.
\Needspace{18\baselineskip}
\begin{theorem}\label{thm:sharp-determinant-power-concentration}\index{Gauduchon duality!critical-power concentration}
On $X=(\mathbb C/(\mathbb Z+\mathrm i\mathbb Z))^n$, with coordinates $z_j=x_j+\mathrm iy_j$, let $\beta_j=\mathrm dx_j\wedge\mathrm dy_j$, $\beta=\sum_j\beta_j$, $V=n!$ and $\mathrm dV=\beta^n/V$. There are K\"ahler forms
$A_\varepsilon=\beta+\ddc w_\varepsilon$,
$0<\varepsilon<1$, such that the currents $T_\varepsilon=A_\varepsilon^{n-1}/V$ belong to
$\mathscr G_\beta$ and satisfy
\[
 \int_X k_{T_\varepsilon}^{n/(n-1)}\,\mathrm dV=V^{-1/(n-1)}.
\]
Each $T_\varepsilon$ attains the supremum in \eqref{eq:constant-mass-determinant-profile}, with $h=\beta$, for some smooth positive
$p_\varepsilon$ with $\|p_\varepsilon\|_{L^n(\mathrm dV)}=1$. The potentials $w_\varepsilon$ are uniformly
bounded and converge uniformly to a continuous
$\beta$-plurisubharmonic function $w_0$. The currents converge weakly to
$T_0=(\beta+\ddc w_0)^{n-1}/V\in\mathscr G_\beta$, but
$\Gamma_0(T_0)=0$. The functions $k_{T_\varepsilon}^{n/(n-1)}$ are not uniformly
integrable. Nevertheless the Monge--Amp\`ere masses $\int_X(\beta+\ddc w_\varepsilon)^n$, $0\leq\varepsilon<1$,
all equal $V$.
\end{theorem}
\begin{proof}
Put $r=1-\varepsilon$ and
\[
 P_r(x)=\frac{1-r^2}{1-2r\cos(2\pi x)+r^2},\qquad
 f_\varepsilon=\varepsilon+(1-\varepsilon)P_r.
\]
The positive periodic Poisson kernel has mean one and Fourier series
\[
 P_r=1+2\sum_{m\geq1}r^m\cos(2\pi mx).
\]
Define
\[
 w_\varepsilon=-\frac{2(1-\varepsilon)}{\pi}
       \sum_{m\geq1}\frac{r^m\cos(2\pi mx_1)}{m^2}.
\]
For fixed $\varepsilon>0$ this is smooth. Differentiation gives
$w_\varepsilon''=4\pi(f_\varepsilon-1)$, and hence
\[
 A_\varepsilon=f_\varepsilon(x_1)\beta_1+
        \sum_{j=2}^n\beta_j=\beta+\ddc w_\varepsilon>0.
\]
It is closed, and $\int_X A_\varepsilon^n=V\int_Xf_\varepsilon\,\mathrm dV=V$.
Since $A_\varepsilon^n=Vf_\varepsilon\,\mathrm dV$, the computation of $k_{T_r}$ in the proof of \cref{thm:constant-mass-entire-determinant-profile}, with $f_\varepsilon^{1/n}$ in place of $r$, gives
\[
 k_{T_\varepsilon}=V^{-1/n}f_\varepsilon^{(n-1)/n},
 \qquad k_{T_\varepsilon}^{n/(n-1)}
       =V^{-1/(n-1)}f_\varepsilon.
\]
Thus the integral of $k_{T_\varepsilon}^{n/(n-1)}$ equals $V^{-1/(n-1)}$, which by \cref{thm:constant-mass-entire-determinant-profile} is its largest value over
\emph{all} currents in $\mathscr G_\beta$. Moreover, $p_\varepsilon=f_\varepsilon^{1/n}$ has $L^n$ norm one,
and $\int_X p_\varepsilon k_{T_\varepsilon}\,\mathrm dV=V^{-1/n}$, so
$T_\varepsilon$ attains the supremum in \eqref{eq:constant-mass-determinant-profile}
for $p=p_\varepsilon$.

Since $\sum_{m\geq1}m^{-2}=\pi^2/6$, the Fourier series gives $|w_\varepsilon|\leq\pi/3$, and, as $\varepsilon\to0$, it converges
uniformly to
\[
 w_0=-\frac2\pi\sum_{m\geq1}\frac{\cos(2\pi mx_1)}{m^2}.
\]
This limit is continuous, and it is $\beta$-plurisubharmonic because
$A_0=\beta+\ddc w_0$ is a positive current, as we now show.
The probability measures $f_\varepsilon(x_1)\,\mathrm dx_1$ converge weakly to the
periodic Dirac mass at $x_1=0$, as follows from their Fourier
coefficients, which equal $1$ for $m=0$ and $(1-\varepsilon)^{|m|+1}$ for $m\neq0$. Thus
\[
 A_0=\delta_{x_1=0}\beta_1+\sum_{j=2}^n\beta_j,
\]
where the Dirac measure is tensored with Lebesgue measure in the
remaining real coordinates. Since $w_0$ is bounded, the powers of $A_0$ are
Bedford--Taylor products. Locally
$\beta_1+\ddc w_0=\ddc\psi$ for a bounded subharmonic function $\psi$ of $z_1$ alone,
so $(\beta_1+\ddc w_0)^2=\ddc(\psi\,\ddc\psi)=0$ by the inductive definition of these products. The form $B=\sum_{j\geq2}\beta_j$ is smooth and closed.
Multilinear expansion therefore gives $A^{n-1}=B^{n-1}+(n-1)C\wedge B^{n-2}$, both for $A=A_\varepsilon$ with $C=f_\varepsilon\beta_1$ and for $A=A_0$ with $C=\beta_1+\ddc w_0$. Since $f_\varepsilon\beta_1\to\beta_1+\ddc w_0$ weakly, we obtain
$T_\varepsilon\to T_0=A_0^{n-1}/V$ weakly. This current is
positive, closed and normalized by $\int_X\beta\wedge T_0=1$.

To compute $\Gamma_0(T_0)$, take for $\nu$ in \cref{lem:gaud-determinant-measure-intrinsic} the trace measure $\beta\wedge T_0$. The coefficient matrix $Q$ of $T_0$ is then a positive multiple of the first matrix below on $\{x_1\neq0\}$, which carries the absolutely continuous part, and of the second on $\{x_1=0\}$, which carries the singular part:
\[
 \operatorname{diag}(1,0,\ldots,0),\qquad
 \operatorname{diag}(0,1,\ldots,1).
\]
Their ranks are respectively
one and $n-1$, both less than $n$. Hence $\det Q=0$ almost everywhere for $\nu$, and
formula~\eqref{eq:intrinsic-gaud-determinant-coordinate} gives
$\Gamma_0(T_0)=0$. This is not the Monge--Amp\`ere measure of
$w_0$, which has total mass $V$ by \cref{thm:mass-invariance-characterization}.

Away from $x_1=0$, $f_\varepsilon\to0$ pointwise, while its integral
against $\mathrm dV$ remains one. Thus the functions $k_{T_\varepsilon}^{n/(n-1)}=V^{-1/(n-1)}f_\varepsilon$ are not
uniformly integrable: almost-everywhere convergence to zero and
uniform integrability would force their fixed positive integrals to
vanish. Since $\Gamma_0(T_0)=0$, the functional $T\mapsto\int_X k_T^{n/(n-1)}\,\mathrm dV$ is not
continuous along $T_\varepsilon\to T_0$, although the potentials converge uniformly. The measures
themselves satisfy $\Gamma_0(T_\varepsilon)\to0$ in total mass:
on the complement of a neighborhood of $\{x_1=0\}$ the densities converge uniformly to zero,
and on the neighborhood H\"older's inequality and the uniform $L^{n/(n-1)}$ bound make their integrals small when the neighborhood is thin. Thus the measures $\Gamma_0(T_\varepsilon)$ converge to $\Gamma_0(T_0)$; only the integrals of the critical powers of their densities fail to converge.
\end{proof}

The function $p_\varepsilon$ for which $T_\varepsilon$ is maximizing varies with $\varepsilon$. The construction concerns the complete profile \eqref{eq:constant-mass-determinant-profile} and the critical power integral; it does not give a sequence of exact maximizers for one fixed smooth density.
